% ============================================================================
% Fichier  : partie5/P05_C13_antiforensique_nuage.tex
% Titre    : Anti-forensique, Nuage et Conteneurs
% Auteur   : MINKA MI NGUIDJOÏ Thierry Emmanuel
% Version  : 1.0 -- novembre 2024
% Corrigés : corriges/Corr_C13.tex
% ============================================================================
% Sources : Carrier (2005) pour les attributs NTFS ; ISO/IEC 27037 et 27042
% (paraphrasées) ; modèle de responsabilité partagée de l'informatique en
% nuage (présenté sans citation) ; journaux d'audit de Kubernetes (présentés
% sans citation).
% Faits nouveaux (bible) : sur SRV-FIC01, outil upd.exe daté de 2017 dans
% $STANDARD_INFORMATION, du 13/03/2025 dans $FILE_NAME et $UsnJrnl ; journal
% de sécurité de SRV-FIC01 vidé le 14/03 à 3 h 52 (1102). SOLITRANS,
% février 2026 : instance du portail en nuage compromise, journaux d'audit
% du fournisseur conservés 90 jours, conteneur supprimé par la reprise
% automatique.
% ============================================================================

\chapter{Anti-forensique, Nuage et Conteneurs}
\label{chap:antiforensique}
\chaptermark{Anti-forensique, nuage, conteneurs}

\epigraph{%
    « Effacer une trace en laisse une autre.\\
    Il faut savoir où la chercher. »%
}{Principe du cours}

\niveauChapitre{Expert}{%
    Maîtrise de la partie~II
}

\begin{boxobjectifs}
\begin{itemize}
    \item Reconnaître les principales techniques d'anti-forensique et les
          traces qu'elles laissent.
    \item Détecter une falsification de dates par la confrontation des
          attributs NTFS et des journaux.
    \item Exploiter les traces d'un journal effacé.
    \item Conduire une acquisition dans un environnement en nuage, selon le
          modèle de responsabilité partagée.
    \item Préserver les traces d'un conteneur avant sa disparition.
    \item Coter \NEK{} ce que l'environnement rend inaccessible.
\end{itemize}
\end{boxobjectifs}

\bigskip

% ============================================================================
\section*{Avant de commencer : un outil de 2017}
% ============================================================================

Yaoundé, 28 mars 2025. Dans l'acquisition à froid de SRV-FIC01, M.~FOTSO
trouve l'outil de copie de l'attaquant~: \texttt{upd.exe}, rangé dans un
dossier de mises à jour. L'explorateur affiche une date de création du
4~juin 2017. Si elle est vraie, l'outil était sur le serveur bien avant
l'attaque, et peut-être installé par quelqu'un de la maison. M.~FOTSO ouvre
l'enregistrement MFT du fichier.

\begin{boxalert}
\textbf{Violation typique.} Prendre pour acquise une date affichée par
l'explorateur ou par un outil, sans la confronter aux autres sources qui
datent le même fichier.

\smallskip
\textbf{Règle~:} une date se recoupe. Quand plusieurs sources indépendantes
datent un même objet, leur désaccord est une constatation, et souvent la
trace d'une manipulation.
\end{boxalert}

\separateur

% ============================================================================
\section{Les Techniques d'Anti-forensique}
\label{sec:c13-techniques}
% ============================================================================

\begin{boxdef}
\textbf{\termecle{Anti-forensique}}~: ensemble des techniques qui visent à
empêcher, retarder ou tromper l'investigation~: effacer, masquer, falsifier
ou rendre illisibles les traces.
\end{boxdef}

\begin{tableaucours}{L{3.2cm}L{4.4cm}L{4.6cm}}{%
    \tableauHeader{Technique} &
    \tableauHeader{Ce qu'elle fait} &
    \tableauHeader{Trace qu'elle laisse}
}
Effacement des journaux & Vide le journal de sécurité ou système &
Événement 1102 ou 104, trou dans la numérotation, sources externes
intactes \\
Falsification de dates & Réécrit les dates d'un fichier & Désaccord entre
\texttt{\$STANDARD\_INFORMATION}, \texttt{\$FILE\_NAME} et
\texttt{\$UsnJrnl} \\
Effacement sécurisé & Écrase le contenu d'un fichier & Nom et dates dans
la MFT et le \texttt{\$UsnJrnl}, outil d'effacement lui-même \\
Chiffrement & Rend le contenu illisible & Le conteneur chiffré, l'outil,
parfois la clé en mémoire \\
Exécution en mémoire, outils du système & N'écrit rien sur le disque,
utilise \texttt{powershell}, \texttt{rundll32} & Mémoire, lignes de
commande, journaux de script s'ils sont activés \\
Dissimulation & Cache des fichiers ou des processus & Écarts entre deux
méthodes de lecture (\texttt{pslist} et \texttt{psscan}) \\
\end{tableaucours}
\vspace{-0.8em}
\begin{center}{\small\itshape Tableau~13.1. Techniques d'anti-forensique et traces résiduelles}\end{center}

Le tableau montre une constante~: chaque technique déplace la trace plus
qu'elle ne la supprime. L'enquêteur cherche donc là où l'attaquant ne
pense pas ou ne peut pas agir~: les journaux centralisés hors de la machine,
la mémoire, les structures internes du système de fichiers, les équipements
réseau.

% ============================================================================
\section{Détecter une Falsification de Dates}
\label{sec:c13-dates}
% ============================================================================

Les outils de falsification réécrivent en général les dates de
\texttt{\$STANDARD\_INFORMATION}, que tout programme peut modifier. Les
dates de \texttt{\$FILE\_NAME} et les entrées du \texttt{\$UsnJrnl} sont plus
difficiles à atteindre. Trois indices orientent l'examen~: une date de
création antérieure à l'installation du système, des fractions de seconde
toutes nulles (les outils écrivent souvent des secondes rondes), et un
désaccord entre les attributs.

\begin{boxscenario}{Opération NYANGA. L'outil upd.exe}
\begin{tableaucours}{L{4.6cm}L{4.0cm}L{3.0cm}}{%
    \tableauHeader{Source} &
    \tableauHeader{Date de création} &
    \tableauHeader{Fractions}
}
\texttt{\$STANDARD\_INFORMATION} & 04/06/2017 09:12:00 & ,0000000 \\
\texttt{\$FILE\_NAME} & 13/03/2025 23:58:41 & ,3318274 \\
\texttt{\$UsnJrnl}, création du fichier & 13/03/2025 23:58:41 & \\
Installation du système (registre) & 17/01/2020 & \\
\end{tableaucours}

\medskip
\textbf{Constatations.} La date de \texttt{\$STANDARD\_INFORMATION} précède de
deux ans et demi l'installation du système, et ses fractions sont nulles.
\texttt{\$FILE\_NAME} et le \texttt{\$UsnJrnl} concordent sur le 13~mars
2025 à 23~h~58.

\textbf{Analyse} (confiance élevée). La date de 2017 a été falsifiée.
L'outil a été déposé le 13~mars à 23~h~58, soit soixante-quatorze minutes
avant le début de l'exfiltration. La piste d'un outil ancien, installé par
un employé, tombe.
\end{boxscenario}

% ============================================================================
\section{Un Journal Vidé}
\label{sec:c13-journal}
% ============================================================================

\begin{boxscenario}{Opération NYANGA. Le journal de SRV-FIC01}
\textbf{Constatations.} Le journal de sécurité de SRV-FIC01 commence par un
événement 1102, le 14~mars à 3~h~52, enregistré sous le compte
\texttt{adm.kamga}~: le journal a été vidé, cinq minutes après la fin de
l'exfiltration. Les journaux du contrôleur de domaine, eux, conservent les
authentifications de \texttt{adm.kamga} vers SRV-FIC01 le 13 à 23~h~55 et
le 14 à 1~h~09.

\textbf{Analyse} (confiance élevée). L'attaquant a effacé les traces locales
de ses sessions, mais pas celles que le contrôleur de domaine garde sur les
authentifications. L'effacement lui-même est une constatation, datée et
attribuée à un compte.
\end{boxscenario}

\begin{boxinfo}
\textbf{La centralisation des journaux.} Un journal copié en temps réel sur
un serveur de collecte, hors de portée des comptes d'administration
ordinaires, survit à l'effacement local. C'est la mesure la plus simple
contre l'effacement des journaux, et NYANGA ne l'avait pas.
\end{boxinfo}

% ============================================================================
\section{L'Enquête dans le Nuage}
\label{sec:c13-nuage}
% ============================================================================

Dans l'informatique en nuage, l'organisation ne possède ni les machines ni
le réseau. Le \termecle{modèle de responsabilité partagée} répartit la
sécurité entre le fournisseur et le client~: plus le service est géré par
le fournisseur, moins le client a de prise, et moins l'enquêteur a de
traces directement accessibles.

\begin{tableaucours}{L{2.4cm}L{3.6cm}L{3.2cm}L{3.0cm}}{%
    \tableauHeader{Modèle} &
    \tableauHeader{Ce que gère le client} &
    \tableauHeader{Traces accessibles} &
    \tableauHeader{Traces chez le fournisseur}
}
Infrastructure (IaaS) & Système, applications, données & Disques des
machines (instantanés), journaux du système & Hyperviseur, réseau physique \\
Plateforme (PaaS) & Applications, données & Journaux de l'application et du
service & Système sous-jacent \\
Logiciel (SaaS) & Données, comptes & Journaux d'audit exposés par le
service & Presque tout le reste \\
\end{tableaucours}
\vspace{-0.8em}
\begin{center}{\small\itshape Tableau~13.2. Traces et responsabilité partagée}\end{center}

L'acquisition y change de forme. On ne débranche pas un serveur~: on prend
un instantané de ses disques par l'interface du fournisseur, on exporte les
journaux d'audit de la console, on copie les journaux du réseau virtuel s'ils
sont activés. La mémoire d'une machine virtuelle ne se capture que par un
outil exécuté à l'intérieur, ou par un service spécialisé du fournisseur.
Les durées de conservation par défaut des journaux d'audit sont souvent
courtes, et les données peuvent se trouver dans un autre pays que
l'organisation, ce qui fait intervenir l'entraide internationale
(Chapitre~\ref{chap:normes-droit}).

\begin{boxscenario}{Opération SOLITRANS. Le portail en février 2026}
\textbf{Constatations.} Le 9~février 2026, le portail des transitaires,
hébergé chez un fournisseur en nuage sous forme de conteneurs, présente des
pages modifiées. L'équipe de M.~NJIKE prend un instantané des disques des
machines virtuelles et exporte les journaux d'audit de la console, conservés
90~jours. Le conteneur compromis, lui, a déjà été remplacé~: le système de
reprise automatique l'a détruit et relancé à partir de l'image d'origine
quand il a cessé de répondre, à 4~h~10.

\textbf{Analyse} (confiance élevée). Le contenu du conteneur compromis est
perdu~; seuls restent les journaux que l'orchestrateur et le fournisseur
ont gardés. La reprise automatique, qui sert la disponibilité, a détruit la
preuve.
\end{boxscenario}

% ============================================================================
\section{Les Conteneurs}
\label{sec:c13-conteneurs}
% ============================================================================

Un conteneur est éphémère par conception~: il naît d'une image, vit
quelques minutes ou quelques jours, et disparaît avec ses fichiers et ses
journaux s'ils n'ont pas été envoyés ailleurs. L'enquête doit donc préparer
la préservation avant l'incident.

\begin{tableaucours}{L{4.0cm}L{8.2cm}}{%
    \tableauHeader{Trace} &
    \tableauHeader{Comment la préserver}
}
Image d'origine & Conserver son empreinte (\textit{digest}) et l'image
elle-même dans le registre \\
Modifications du conteneur & Copier la couche modifiable avant de détruire
le conteneur, ou le suspendre au lieu de le redémarrer \\
Journaux de l'application & Les envoyer en continu vers un collecteur
externe \\
Journaux de l'orchestrateur & Activer et conserver le journal d'audit
(créations, suppressions, accès aux secrets) \\
Réseau & Journaux des flux du réseau virtuel \\
\end{tableaucours}
\vspace{-0.8em}
\begin{center}{\small\itshape Tableau~13.3. Préserver les traces d'un conteneur}\end{center}

\begin{boxlecture}{Ce que l'enquête peut encore savoir, NYANGA et SOLITRANS}
\begin{tableaucours}{L{7.4cm}c}{%
    \tableauHeader{Proposition} &
    \tableauHeader{Cote}
}
\texttt{upd.exe} a été déposé sur SRV-FIC01 le 13~mars 2025 & \OUI{} \\
\texttt{upd.exe} était présent en 2017 & \NON{} \\
Le journal de SRV-FIC01 a été vidé par le compte \texttt{adm.kamga} &
\OUI{} \\
Les sessions de l'attaquant sur SRV-FIC01 sont datées & \OUI{} \\
Le conteneur du portail de SOLITRANS contenait un outil de l'attaquant &
\NEK{} \\
Le point d'entrée sur le portail de SOLITRANS est identifié & \NEK{} \\
\end{tableaucours}

\medskip
\textbf{Lecture.} Sur NYANGA, l'anti-forensique a échoué~: chaque trace
effacée ou falsifiée a été recoupée ailleurs. Sur SOLITRANS, deux \NEK{}
tiennent à l'architecture elle-même~: ce que la reprise automatique a
détruit ne se reconstitue pas.
\end{boxlecture}

% ============================================================================
\section{Ce que ces Défis demandent à chaque Métier}
\label{sec:c13-metiers}
% ============================================================================

\begin{boxmetier}{Analyste en réponse à incident}
\textbf{Face à l'anti-forensique}
\begin{enumerate}
    \item Recouper chaque date par au moins deux sources indépendantes.
    \item Chercher les journaux hors de la machine (contrôleur, collecteur,
          réseau).
    \item Traiter chaque effacement comme une constatation datée.
    \item Dans le nuage, exporter d'abord les journaux d'audit, dont la
          conservation est courte.
    \item Suspendre un conteneur suspect au lieu de le laisser redémarrer.
\end{enumerate}
\end{boxmetier}

\begin{boxmetier}{RSSI}
Centraliser les journaux hors de portée des comptes d'administration~;
négocier avec les fournisseurs en nuage la durée de conservation des
journaux d'audit~; régler la reprise automatique des conteneurs pour
qu'elle isole avant de détruire.
\end{boxmetier}

\begin{boxmetier}{Application OPJ/Expert Judiciaire}
Pour les données détenues par un fournisseur à l'étranger, l'expert indique
dans son rapport ce qu'il n'a pas pu examiner et pourquoi, et l'OPJ engage
sans délai la demande de conservation puis l'entraide.
\end{boxmetier}

\separateur

\begin{boxresume}
\textbf{Ce qu'il faut retenir de ce chapitre~:}
\begin{itemize}
    \item L'anti-forensique déplace les traces plus qu'elle ne les
          supprime~; on les cherche là où l'attaquant n'agit pas.
    \item Une falsification de dates se détecte par le désaccord entre
          \texttt{\$STANDARD\_INFORMATION}, \texttt{\$FILE\_NAME} et
          \texttt{\$UsnJrnl}.
    \item Un journal vidé laisse un événement daté et attribué~; les
          journaux centralisés survivent.
    \item Dans le nuage, l'acquisition passe par les instantanés et les
          journaux d'audit, aux durées courtes.
    \item Un conteneur se préserve avant l'incident~; la reprise automatique
          peut détruire la preuve.
\end{itemize}
\end{boxresume}

\begin{boxinfo}
\textbf{Points de contrôle.} Vous avez compris ce chapitre si vous savez~:
\begin{itemize}
    \item démontrer une falsification de dates à partir de trois sources~;
    \item dire ce qu'un client peut acquérir selon le modèle de service~;
    \item préparer la préservation d'un environnement de conteneurs.
\end{itemize}
\end{boxinfo}

% ============================================================================
\section*{Exercices}
% ============================================================================

\begin{boxexo}[title={\ding{45}\quad Exercice~13.1. Trace résiduelle \niveaubase}]
Pour chaque technique, citez une trace résiduelle à chercher~: (a) un
journal de sécurité vidé~; (b) un fichier écrasé par un outil d'effacement~;
(c) un script exécuté en mémoire par \texttt{powershell}~; (d) une archive
chiffrée.
\end{boxexo}

\begin{boxexo}[title={\ding{45}\quad Exercice~13.2. Vraie ou fausse date \niveaumoyen}]
Un fichier \texttt{rapport.xlsx} présente une création au 12/03/2025
14:05:33,8812094 dans \texttt{\$STANDARD\_INFORMATION} et au 12/03/2025
14:05:33,8812094 dans \texttt{\$FILE\_NAME}~; le \texttt{\$UsnJrnl} ne
contient plus d'entrée aussi ancienne.
\begin{enumerate}[(a)]
    \item La date est-elle falsifiée~?
    \item Pourquoi le \texttt{\$UsnJrnl} ne permet-il pas de trancher~?
    \item Comment cotez-vous la proposition « le fichier a été créé le
          12~mars »~?
\end{enumerate}
\end{boxexo}

\begin{boxexo}[title={\ding{45}\quad Exercice~13.3. Acquérir dans le nuage \niveaumoyen}]
Une machine virtuelle de SOLITRANS, hébergée en IaaS, est suspectée.
Décrivez, dans l'ordre, les acquisitions à faire, et dites pour chacune ce
qui serait perdu si on la faisait trop tard.
\end{boxexo}

\begin{boxexo}[title={\ding{45}\quad Exercice~13.4. Préparer les conteneurs \niveauexpert}]
\begin{enumerate}[(a)]
    \item Rédigez cinq règles de préparation qui auraient permis d'examiner
          le conteneur compromis du portail de SOLITRANS.
    \item La reprise automatique sert la disponibilité. Proposez un réglage
          qui concilie disponibilité et préservation.
    \item Reprenez l'encadré « Ce que l'enquête peut encore savoir » et dites
          lesquelles des deux cotes \NEK{} de SOLITRANS vos règles auraient
          pu changer.
\end{enumerate}
\end{boxexo}

\bigskip

\begin{boxinfo}
\textbf{Pour aller plus loin.}
\begin{itemize}
    \item B.~Carrier, \textit{File System Forensic
          Analysis}~\cite{carrier2005}, sur les attributs NTFS.
    \item L'ISO/IEC~27042~\cite{iso27042_2015}, sur l'analyse et
          l'interprétation.
\end{itemize}
\end{boxinfo}

\begin{boxinfo}
\textbf{Transition.} Tous les outils sont en place~: acquérir, analyser,
dater, remonter à la faille, contourner l'anti-forensique. Le
Chapitre~\ref{chap:enquete} les rassemble dans une enquête complète, du
premier appel au dernier scellé.
\end{boxinfo}
